\chapter{Introduction}

\label{chapter: Introduction}
Altermagnetism has become an emerging topic in the field of spintronics and magnetism. Altermagnets combine properties of ferro- and antiferromagnets: They have spin-split bands, while their magnetic order is fully compensated. Because of their vanishing or small net magnetization, they represent promising candidates for spintronics due to their weak magnetic stray fields and therefore are robust against external magnetic fields.\\
Despite their compensated magnetic order, they can exhibit properties that are common in ferromagnets. One example for this is the anomalous Hall effect (AHE), which was studied for different compounds in Refs.~\cite{PhysRevB.111.184407, Quelle2, MnPtSOC}. \\
Altermagnetism has also attracted interest in the context of superconductivity \cite{PhysRevB.108.224421},  where the superconducting diode effect (SDE)  was studied in altermagnets \cite{cv8s-tk4c, Quelle3}. \\

Originally, altermagnetism was introduced for collinear systems \cite{PhysRevX.12.040501} and was later extended to noncollinear magnetic systems \cite{Cheong2024, Cheong2025}. Due to their geometry, kagome lattices support $120^\circ$ magnetic order and therefore supply a great platform to achieve noncollinear altermagnets. Depending on their magnetic order and the corresponding symmetries, they can realize different types of altermagnetism, which are classified in Ref. \cite{Cheong2025}. A few examples of noncollinear altermagnets include \ce{Mn3Pt}, \ce{Mn3Ge} and \ce{Mn3Sn}.

For the study of the compounds, different theoretical methods can be used, including first-principles calculations \cite{Quelle2, Yang_2017}, Monte Carlo simulations \cite{9d8q-cjk7} and tight-binding models \cite{PhysRevB.99.165141}.
In this thesis, a tight-binding model is used to provide a simple description of the electronic structure, to simplify symmetry analysis, and to allow the effects of spin-orbit coupling and different magnetic orders to be studied separately. \\

In this thesis, we aim to describe the electronic-structure properties, namely the electronic band structure, Fermi-surface spin texture and Berry curvature, for the magnetic textures of the T1 and T2 phases of \ce{Mn3Pt} and of \ce{Mn3Ge} and \ce{Mn3Sn}, on regular and breathing kagome lattices with and without Rashba-SOC. The model for \ce{Mn3Pt} follows Ref. \cite{Quelle3}, which serves as a guide in this thesis. 
In Chapter \ref{chapter:Theory}, the tight-binding model and the Berry curvature are introduced, together with the different types of altermagnetism and their relevant symmetries.
Then, in Chapter \ref{chapter:Results}, the model introduced above is applied to the different magnetic textures, the relevant symmetries are identified and
their predictions are compared with the numerical results.
Finally, Chapter \ref{chapter: conclusion} summarizes the main findings and gives an outlook on further extensions.\\






